\begin{tabular}{lp{.25\textwidth}p{.18\textwidth}p{.4\textwidth}r}
\toprule
& Definition & Sample & Challenges \& Limitations & N\tn[\dag] \\
% SM: in text above, we talk about limitations.
%     do we mean limitations or challenges or both?
%     and if both, should we separate out what is a challenge vs. what is a limitation in the table?
% JK: Great point that these are not the same (challenges vs limitations)!
%     However, I think for our purposes here, it is OK to combine them as "challenges & limitations"?
%     I would be worried about having readers a bit distracted by the distinction?
\midrule
1. & mean total duration
   & former sex workers
   & \tabullet very difficult to sample this population
   & 0 \\
\midrule
2. & inverse of the proportion who started selling sex within the past year
   & active sex workers
   & \tabullet this proportion is rarely reported
     \tabullet sensitive to undersampling early sex work
     \tabullet not adjusted for length-biased sampling\tn[\ddag]
   & 16 \\
\midrule
3. & double the mean/median (right-censored) duration
   & active sex workers
   & \tabullet median $<$ mean for skewed distributions
% SM: so the limitation only apply to when using double the median?
% JK: yes
     \tabullet not adjusted for length-biased sampling\tn[\ddag]
   & 93 \\
\bottomrule
\end{tabular}
\smallskip\floatfoot{
\par\tnd[\dag]{%
  Number of estimates of duration selling sex synthesized by \cite{Fazito2012};
  corresponding numbers were: 0, 1, 8 for duration buying sex.}
\par\tnd[\ddag]{%
  More likely to sample individuals with longer durations \cite{Kafadar2009}.}}
